\section{总结与延伸}

\subsection{核心观点回顾}

Birgitta B\"{o}ckeler 对 OpenAI harness engineering 实践的分析揭示了 AI 辅助编程的多个深层趋势：

\begin{enumerate}[leftmargin=1.5em]
    \item \textbf{Harness 是系统性的}：由上下文工程、架构约束和``垃圾回收''三大组件构成，结合确定性和 LLM 驱动的方法
    \item \textbf{约束产生信任}：放弃``生成任意代码''的灵活性，换取可维护性和可靠性
    \item \textbf{技术栈将收敛}：AI 友好性可能成为技术选型的重要标准
    \item \textbf{两个世界的分化}：新建应用可以 harness-native，遗留系统面临改造挑战
    \item \textbf{严谨性在迁移}：工程严谨性从编码层迁移到环境和控制系统设计层
    \item \textbf{功能验证仍有缺口}：当前的 harness 实践偏重内部质量，外部质量验证有待加强
\end{enumerate}

\subsection{对实践者的启示}

\begin{warningbox}{保持批判性思维}
OpenAI 对自身产品的成功案例描述需要保持审慎态度------作为 Codex 的开发者，他们在让我们相信 AI 可维护代码方面有天然的利益动机（vested interest）。在参考其实践时，应关注方法论本身的合理性，而非仅仅依赖其宣称的成功指标。
\end{warningbox}

对于希望开始 harness engineering 实践的团队，以下是建议的优先级：

\begin{enumerate}[leftmargin=1.5em]
    \item \textbf{评估现状}：盘点现有的 linter、pre-commit hook、CI 检查等
    \item \textbf{显式化架构约定}：将隐式的团队约定转化为可自动化的规则
    \item \textbf{引入结构性测试}：使用 ArchUnit 等工具验证架构边界
    \item \textbf{建设知识库}：维护代码库中的设计决策记录和架构文档
    \item \textbf{迭代改进}：将 AI agent 的失败视为 harness 改进的信号
\end{enumerate}

\subsection{拓展阅读}

\begin{itemize}[leftmargin=1.5em]
    \item OpenAI: ``Harness engineering: leveraging Codex in an agent-first world''
    \item Mitchell Hashimoto: Harness 概念的早期讨论
    \item Chad Fowler: ``Relocating Rigor''------严谨性迁移的思考
    \item Martin Fowler 博客 ``Exploring Gen AI'' 系列
    \item ArchUnit 官方文档------结构性测试框架
    \item Birgitta B\"{o}ckeler: \href{https://martinfowler.com/articles/exploring-gen-ai/harness-engineering.html}{原文链接}
\end{itemize}

%% --- Fin del contenido --- %%

\end{document}
